
%
% $Log: layout.tex,v $
% Revision 1.4  2002/10/21 15:26:44  mjm
% to ohio, revised
%
% Revision 1.3  2001/08/15 18:07:03  mjm
% APPM preprint 466
%
% Revision 1.2  1999/08/09 19:55:27  mjm
% beta release
%
% Revision 1.1  1999/07/12 22:04:36  mjm
% Initial revision
%
%
%

%%%%%%%%%%%%%%%%%%%%%%%%%%%%%%%%%%%%%%%%%%%%%%%%%%%%%%%%%%%%%%%%%%%%%%%%%
%\subsection*{Laying out the Problem}

\handout{Laying out the Problem}

Presentation is very important in mathematics, just as it is in other
fields. The main reason for this is clarity. Good presentation allows you
to  communicate more clearly the content of your work.
A secondary reason is the appearance of competence. Careful presentation
makes the reader think you were also careful with the content of your
work. Although good presentation is rarely successful at bluffing past poor
content, poor presentation can easily ruin good content.

The first thing the reader evaluates is the overall visual layout of the
problem. Exactly how polished this should be will depend on the purpose of
your writing. For our purposes, we have the following rules:
\begin{itemize}
	\item
Put the good problem on its own sheet of paper, separate from any
other problems. Regular quality paper is fine, 
but there should be no ragged edges or tears.
	\item
If you need more than one page, staple (no folded corners!) them together
and put your name and the page number on each page.
	\item
Leave margins on all sides.
	\item
Print neatly or type. Do not switch colors or from pen to pencil in the
middle of the problem. (You can use different colors to highlight if you wish.)
	\item
If you had to cross out material or erased a lot and left smudges, rewrite
the problem. (It is a good habit to solve the problem first on scrap paper
and then copy it neatly.)
\end{itemize}
%
The reader next needs to know what it is they are reading. On top of the first
page, put:
\begin{enumerate}
	\item
your full name,
	\item
the course and section or recitation number, and
	\item
the assignment number and/or the date the assignment is due.
\end{enumerate}
%
Next is the format for the problem itself:
\begin{itemize}
	\item
Label the problem with the chapter, section, and problem number.
	\item
Write out the entire question, including any instructions. If the question
refers to another problem, include the relevant information from that
problem. The goal is to make your work as self-contained as possible, so
the reader does not need to look anything up. 
	\item
Do the problem in some logical order. Do not do the problem in several
disjoint pieces connected by arrows.
\end{itemize}

All these rules may seem picky. Once you have learned this way to lay out
your work, you will understand the principles behind these rules. Then you
will know when to change the rules.

{\bf Acknowledgements:} If you had help with this problem, or in any
way include work that is not your own, then give proper credit. Not
only is this the nice thing to do, but also it will protect you from
allegations of cheating.

Good presentation is important in mathematics. It lets you communicate your
work more clearly and lends it credibility.


%%%%%%%%%%%%%%%%%%%%%%%%%%%%%%%%%%%%%%%%%%%%%%%%%%%%%%%%%%%%%%%%%%%%%%%%%
\secondpage

\subsection*{Example:}
%{\bf Example:}

\begin{goodexample}

\begin{flushright}
Full Name\\
Course and section or recitation number\\
Assignment number\\
Date\\
\end{flushright}

\noindent{\bf Section number, problem number.}
 Full text of the problem \ldots blah \ldots blah \ldots blah \ldots blah.

\vspace{\baselineskip}

In this problem we \ldots blah \ldots blah \ldots blah \ldots blah \ldots
blah \ldots blah \ldots blah \ldots blah. \\ 
\fbox{\rm There is a separate handout on introductions.}

\vspace{\baselineskip}

Since \ldots blah \ldots blah \ldots we know \ldots blah \ldots blah.\\
\fbox{\rm There is a separate handout on logical connectives.}

\begin{align*}
\implies & \mbox{blah}\\
\implies & \mbox{blah blah}\\
\therefore & \quad \mbox{blah!}\\
& \fbox{\rm There is a separate handout on mathematical symbols.}
\end{align*}
From this we can conclude \ldots blah \ldots blah \ldots blah \ldots blah
and so $x=3\pi\ldots$ \ldots blah. \\
\noindent\fbox{\parbox{\textwidth}{\rm
There is a separate handout on incorporating mathematics into sentences and other issues of `flow'.}}

\vspace{\baselineskip}

By graphing our solution, we can see that \ldots blah.
\begin{center}
\setlength{\unitlength}{0.5mm}
\begin{picture}(100,100)(0,0)
	\put(0,0){\line(1,0){100}}
	\put(0,0){\line(0,1){100}}
	\put(100,0){\line(0,1){100}}
	\put(0,100){\line(1,0){100}}
	\put(8,60){\mbox{\rm There is a separate handout}}
	\put(20,40){\mbox{\rm on graphing.}}
\end{picture}
\end{center}


\vspace{\baselineskip}

By using \ldots blah \ldots blah \ldots we were able to solve this
problem \ldots blah \ldots blah \ldots blah \ldots blah.\\ 
\fbox{\rm There is a separate handout on conclusions.}

\vspace{\baselineskip}
{\bf Acknowledgements:} I worked with so-and-so on this problem and also had 
help from so-and-so.

\end{goodexample}


